\begin{proof}
Notation as in the statement; $s = Q(w)$, $|w| = \sqrt{s}$, and $|S^{3}| = 2\pi^{2}$ is the area of the unit sphere of $\R^{4}$. Identities between kernels off the origin are identities of finite sums of polynomials in the components times $s^{a}(\ln s)^{k}$ and are computed exactly in the check; pairings with the test functions of the check are exact rational multiples of $\pi^{2}$. The photon propagator in a covariant gauge is used inside the loop under the Krein-form premise \claimref{synthesis.krein-form-premise-indefinite-inner-product-for-gauge-fields}; no step below uses an inner product, so the premise bears on the reading of the kernels as electrodynamics and not on the identities.

\emph{Step 1 (clause (i)).} In $\Cl(4,0)$, $\sum_{a}\gamma_{a}w\gamma_{a} = -2w$, $\sum_{a}\gamma_{a}\gamma_{a} = 4$ and $www = Q(w)w$. Hence $\sum_{a,b}\gamma_{a}w\gamma_{b}\big(\tfrac{1+\xi}{2}\delta_{ab} + (1-\xi)w_{a}w_{b}/s\big) = -(1+\xi)w + (1-\xi)w = -2\xi w$ and $\sum_{a,b}\gamma_{a}\gamma_{b}\big(\tfrac{1+\xi}{2}\delta_{ab} + (1-\xi)w_{a}w_{b}/s\big) = 2(1+\xi) + (1-\xi) = 3+\xi$. With the propagator of the Dirac field of Proposition \claimref{field-theory.second-quantisation}, the vertex of Proposition \claimref{field-theory.coupled-system-tree-level}, and the expansion of the propagator used in Theorem \claimref{synthesis.one-loop-vacuum-polarisation-extends-without-a-regulator-and-ward-preserving-extensions-differ-only-by-the-renormalisation-constants}(ii), $S = \big(\frac{1}{2\pi^{2}s^{2}} - \frac{m^{2}}{8\pi^{2}s}\big)w + \frac{m}{4\pi^{2}s} + (\text{scaling degree at most } 0)$, this gives the two leading terms of $\Sigma$ with coefficients $\frac{1}{2\pi^{2}}\cdot\frac{1}{4\pi^{2}}\cdot(-2\xi) = -\frac{\xi}{4\pi^{4}}$ and $\frac{1}{4\pi^{2}}\cdot\frac{1}{4\pi^{2}}\cdot(3+\xi)$; the rest has scaling degree at most $3$. The scaling degrees $3$ of $S$ and $2$ of $D$ add to $5$ for $\Sigma$ and to $8$ for $\Lambda$ at the origin of $\R^{8}$; on $\R^{8}\setminus\{0\}$ at most one of $u$, $v$, $u+v$ vanishes at a point, and each factor is locally integrable in $\R^{4}$ (degrees $3$, $3$, $2$, all below $4$), so $\Lambda$ is locally integrable there. For $m = 0$ and $\xi = 0$ the first identity gives $\Sigma = 0$ off the origin.

\emph{Step 2 (the divergence off the origin).} $S$ is locally integrable with $\nabla S = \delta - mS$ and $S\overleftarrow{\nabla} = \sum_{a}\partial_{a}S\,\gamma_{a} = \delta - mS$ as distributions (for $m = 0$: $w/s^{2}$ is monogenic off the origin and its flux through a sphere is $|S^{3}|$, computed). Near a point of $\R^{8}\setminus\{0\}$ with $v = 0$ the factors $S(u)$ and $D_{u+v}$ are smooth, and similarly at $u = 0$; where neither vanishes all derivatives are classical. So, the mass terms cancelling,
\[
\sum_{a}(\partial_{v_{a}} - \partial_{u_{a}})\,S(u)\gamma_{a}S(v) = S(u)\,\nabla S(v) - \big(S\overleftarrow{\nabla}\big)(u)\,S(v) = \delta(v)S(u) - \delta(u)S(v),
\]
and multiplying by the photon kernel and summing over $a$, $b$ gives $\operatorname{div}\Lambda = \delta(v)\Sigma(u) - \delta(u)\Sigma(v)$ on $\R^{8}\setminus\{0\}$. Only the dependence of the photon kernel on $w = u + v$ alone was used, since $\partial_{v_{a}} - \partial_{u_{a}}$ annihilates functions of $u + v$.

\emph{Step 3 (the zero-transfer integral).} Fix $n$ and put $G(w) = \int S(x-z)\,n\,S(z-y)\,d^{4}z - (n\cdot w)S(w)$, $w = x - y$. The integral converges absolutely for $w \neq 0$: the integrand is $O(|z|^{-6})$ at infinity (exponentially small for $m > 0$) and has integrable singularities of degree $3$ at $z = x$ and $z = y$; it is locally integrable in $w$, of size $O(|w|^{-2})$ near $0$. It is the convolution of $S$ with $nS$, so $(\nabla + m)$ applied to it gives $nS$. Also $(\nabla + m)\big((n\cdot w)S\big) = nS + (n\cdot w)(\nabla + m)S = nS + (n\cdot w)\delta = nS$. Hence $(\nabla + m)G = 0$ on $\R^{4}$ as a distribution, so $(-\nabla^{2} + m^{2})G = (-\nabla + m)(\nabla + m)G = 0$. For $m > 0$ a tempered solution of this equation is zero. For $m = 0$, $G$ is a tempered harmonic distribution, hence a polynomial, and it is homogeneous of degree $-2$, hence zero. So $\int S(x-z)\,n\,S(z-y)\,d^{4}z = (n\cdot w)S(w)$ (the position-space form of the derivative of the propagator along $n$ in momentum). The photon kernel depends on $w$ alone, so multiplying by it and summing gives the displayed formula of (ii), with $z = \frac{x+y}{2} - t$.

\emph{Step 4 (the extensions).} $\bar\Sigma^{\ell}$: near the origin the bracket is $O(|w|^{2})$ against a kernel of degree $5$ in $\R^{4}$, integrable; elsewhere the kernel is locally integrable and $g$ has compact support. $\bar\Lambda^{\ell}$: near the origin the bracket is $O(|u| + |v|)$ against a kernel of degree $8$ in $\R^{8}$, integrable; on the cylinder $Q(u+v) < \ell^{2}$ far from the origin the integrand is $O(|t|^{-6})$ times the locally integrable $|w|^{-2}$, integrable over $\R^{4}_{t}\times\{|w| < \ell\}$. On test functions vanishing near the origin the subtracted terms vanish, so both are extensions. Each differs from the extension with a smooth compactly supported cut-off function $\chi$ equal to $1$ near the origin (the construction of Fact \claimref{synthesis.imported-epstein-glaser-extension-theorem-and-ward-preserving-time-ordered-products}(ii)) by the pairing of the kernel with $(\chi - \theta)$ times a polynomial of degree at most $1$, respectively $0$, a convergent integral that depends on the test function only through its Taylor coefficients at the origin of order at most $1$, respectively $0$: a combination of $\delta$ and $\nabla\delta$, respectively $\delta(u)\delta(v)$. So the scaling degrees are those of the extensions of that Fact. The cut-off regions are invariant under the orthogonal group and the kernels are covariant, so the extensions are covariant; the oddness of $w \mapsto w/s^{3}$ makes the term $g(0)$ of the subtraction drop out of the leading part of $\bar\Sigma^{\ell}$.

\emph{Step 5 (clause (iii)).} Use the coordinates $w = u + v$, $t = (u - v)/2$ on $\R^{8}$ (unit Jacobian), in which $\partial_{u_{a}} - \partial_{v_{a}} = \partial_{t_{a}}$. Let $f$ be a test function and $f_{a} = \partial_{t_{a}}f$. By definition of the distributional derivative,
\[
\langle\operatorname{div}\bar\Lambda^{\ell},f\rangle = \sum_{a}\langle\bar\Lambda^{\ell}(\gamma_{a}),f_{a}\rangle = \sum_{a}\int\Lambda(\gamma_{a};u,v)\big(f_{a}(u,v) - \theta(\ell^{2} - Q(w))f_{a}(0,0)\big)\,d^{4}w\,d^{4}t .
\]
The integrand is absolutely integrable (Step 4), so by Fubini's theorem the $t$-integral may be done first, for each $w \neq 0$. There each of the two terms converges separately. For the first, $\Lambda(\gamma_{a};\cdot)$ at fixed $w \neq 0$ is a locally integrable function of $t$, decaying as $|t|^{-6}$, whose distributional divergence in $t$ is, by the computation of Step 2 done in $t$ at fixed $w$ (the two singular points $t = w/2$ and $t = -w/2$ are distinct for $w \neq 0$ and the other factors are smooth near each; this is not a restriction of a distribution on $\R^{8}$), $-\big(\delta(\tfrac{w}{2}-t) - \delta(\tfrac{w}{2}+t)\big)\Sigma(w)$; so
\[
\sum_{a}\int\Lambda(\gamma_{a};u,v)\,\partial_{t_{a}}f\,d^{4}t = \Sigma(w)\big(f|_{t = w/2} - f|_{t = -w/2}\big) = \Sigma(w)\,\hat f(w), \qquad \hat f(w) := f(u{=}w,v{=}0) - f(u{=}0,v{=}w).
\]
For the second, Step 3 gives $\sum_{a}f_{a}(0,0)\int\Lambda(\gamma_{a};u,v)\,d^{4}t = \sum_{a}f_{a}(0,0)\,w_{a}\,\Sigma(w)$. Now $\hat f$ is a test function on $\R^{4}$ with $\hat f(0) = 0$ and, by the chain rule, $\partial_{a}\hat f(0) = (\partial_{w_{a}} + \tfrac12\partial_{t_{a}})f(0,0) - (\partial_{w_{a}} - \tfrac12\partial_{t_{a}})f(0,0) = f_{a}(0,0)$. Therefore
\[
\langle\operatorname{div}\bar\Lambda^{\ell},f\rangle = \int\Sigma(w)\big(\hat f(w) - \theta(\ell^{2} - Q(w))(\hat f(0) + w\cdot\nabla\hat f(0))\big)\,d^{4}w = \langle\bar\Sigma^{\ell},\hat f\rangle ,
\]
and $\langle\bar\Sigma^{\ell},\hat f\rangle$ is by definition the pairing of $\delta(v)\bar\Sigma^{\ell}(u) - \delta(u)\bar\Sigma^{\ell}(v)$ with $f$. This is clause (iii). No property of $\xi$, of $m$ or of the form of the photon kernel beyond Steps 2 and 3 was used. The identity is built in by the two regions: the subtraction region of the vertex is the product $\{|w| < \ell\}\times\R^{4}_{t}$, so the integration by parts in $t$ at fixed $w$ meets no boundary (a ball in $\R^{8}$ would leave a boundary term), and by Step 3 the subtracted vertex term integrates to the Taylor subtraction of the self-energy on the ball of the same radius (with different radii the shell between them remains, Step 6). In terms of the stored constants the choice is $c_{V} = -c_{2}$.

\emph{Step 6 (clause (iv)).} For a kernel $h$ homogeneous of degree $-4 - j$ and the subtraction of the Taylor polynomial $T_{j}g$ of order $j$ on the ball of radius $\ell = 1/\mu$, $\mu\,\partial_{\mu}\langle[h]_{\ell},g\rangle = \ell\,\partial_{\ell}\int_{|w| < \ell}h\,T_{j}g = \int_{|w| = 1}h\,T_{j}g\,d\Omega$. The average of $w_{a}w_{b}$ over the unit sphere is $\delta_{ab}/4$. So $\mu\,\partial_{\mu}\langle[w_{a}/s^{3}]_{\ell},g\rangle = \frac{|S^{3}|}{4}\partial_{a}g(0) = -\frac{\pi^{2}}{2}\langle\partial_{a}\delta,g\rangle$ and $\mu\,\partial_{\mu}\langle[n\cdot w\,w/s^{3}]_{\ell},g\rangle = \frac{\pi^{2}}{2}\,n\,g(0)$, and with the coefficient $-\xi e^{2}/(4\pi^{4})$: $\mu\,\partial_{\mu}\bar\Sigma^{\ell} = \frac{\xi e^{2}}{8\pi^{2}}\nabla\delta$. For the vertex, $\bar\Lambda^{\ell}(n) - \bar\Lambda^{\ell'}(n)$ paired with $f$ is $-f(0,0)$ times the integral of $\Lambda$ over $\ell' \le |w| < \ell$ and all $t$, which by Step 3 is the integral of $(n\cdot w)\Sigma(w)$ over the shell: $\mu\,\partial_{\mu}\bar\Lambda^{\ell}(n) = -\frac{\xi e^{2}}{4\pi^{4}}\cdot\frac{\pi^{2}}{2}\,n\,\delta(u)\delta(v)$. The two are opposite, and inserting $\bar\Lambda^{\ell_{V}} = \bar\Lambda^{\ell_{\Sigma}} + \frac{\xi e^{2}}{8\pi^{2}}\ln(\ell_{V}/\ell_{\Sigma})\,n\,\delta(u)\delta(v)$ into (iii) gives the stated defect, which is a nonzero distribution unless its coefficient vanishes. Under $\partial \mapsto$ momentum the self-energy has $\ln\mu^{2}$ with coefficient $\xi e^{2}/(16\pi^{2})$ times the momentum vector. The mass term $\frac{(3+\xi)e^{2}m}{16\pi^{4}}s^{-2}$ has $\mu\,\partial_{\mu} = \frac{(3+\xi)e^{2}m}{16\pi^{4}}|S^{3}|\delta$; the remaining terms of $\Sigma$ for $m > 0$ have scaling degree at most $3$ and contribute to $\mu\,\partial_{\mu}$ terms of relative order $m^{2}\ell^{2}$ (up to logarithms).

\emph{Step 7 (clause (v)).} Take $\ell = 1$; the general case follows by scaling. (a) The distributions $[1/s^{2}]_{1}$ and $-\frac14\nabla^{2}(\ln s/s)$ are rotation-invariant extensions of $1/s^{2}$ (the second by the identity $\nabla^{2}(\ln s/s) = -4/s^{2}$ off the origin, computed), so they differ by $c\,\delta$; pairing with one test function with $g(0) \neq 0$ gives $c = \pi^{2}$ (computed, two test functions). Since $\nabla^{2}(1/s) = -4\pi^{2}\delta$, replacing $\ln s$ by $\ln(M^{2}s)$ adds $\pi^{2}\ln M^{2}\,\delta$: equality at $\ln M^{2} = 1$. (b) $[w/s^{3}]_{1}$ and $\frac{1}{16}\nabla\nabla^{2}(\ln s/s)$ are covariant extensions of $w/s^{3}$ (identity computed), odd under $w \mapsto -w$, so they differ by $c'\,\nabla\delta$; pairing gives $c' = -\frac{3\pi^{2}}{8}$ (computed), and $\ln M^{2}$ adds $-\frac{\pi^{2}}{4}\ln M^{2}\,\nabla\delta$: equality at $\ln M^{2} = \frac32$. (c) $w_{a}w_{b}/s^{3} = (w_{a}w_{b} - \frac14\delta_{ab}s)/s^{3} + \frac14\delta_{ab}/s^{2}$; the first term equals $\frac18(\partial_{a}\partial_{b} - \frac14\delta_{ab}\nabla^{2})(1/s)$ off the origin (computed), and the distributional derivative of the locally integrable homogeneous $1/s$ is a homogeneous extension of it; its cut-off extension does not depend on $\ell$ (zero sphere average) and the difference of the two is a traceless symmetric covariant multiple of $\delta$, hence zero (confirmed by exact pairing). So $[w_{a}w_{b}/s^{3}]_{1}$ is that homogeneous distribution plus $\frac14\delta_{ab}[1/s^{2}]_{1}$, with the scale of (a). (d) By Steps 3 and 5 the zero-transfer vertex $\int\bar\Lambda^{\ell}(n)\,d^{4}t$ is $-\frac{\xi e^{2}}{4\pi^{4}}[(n\cdot w)w/s^{3}]_{\ell}$, and $(n\cdot w)\bar\Sigma^{\ell}$ is the same distribution (their defining integrals coincide). In terms of the scales: $\ln(M_{\Sigma}^{2}\ell^{2}) = \frac32$, $\ln(M_{V}^{2}\ell^{2}) = 1$, so $\ln(M_{\Sigma}^{2}/M_{V}^{2}) = \frac12$. Independently of any cut-off: by the Leibniz rule for a polynomial times a distribution, $w_{b}\,\partial_{a}\nabla^{2}F = \partial_{a}\big(\nabla^{2}(w_{b}F) - 2\partial_{b}F\big) - \delta_{ab}\nabla^{2}F$ with $F = \ln(M^{2}s)/s$, and the check pairs both sides of
\[
\tfrac{1}{16}\,w_{b}\,\partial_{a}\nabla^{2}\frac{\ln(M_{\Sigma}^{2}s)}{s} = \tfrac18\big(\partial_{a}\partial_{b} - \tfrac14\delta_{ab}\nabla^{2}\big)\frac1s - \tfrac{1}{16}\delta_{ab}\nabla^{2}\frac{\ln(M_{V}^{2}s)}{s}
\]
exactly: the two sides differ by $\frac{\pi^{2}}{4}\big(\ln(M_{\Sigma}^{2}/M_{V}^{2}) - \frac12\big)\delta_{ab}\delta$, zero exactly at $\frac12$. If the two scales are not so related, the zero-transfer form of the identity fails by a multiple of $\xi\,n\,\delta$, hence so does (iii); for $\xi \neq 0$ the relation is necessary and sufficient within these two one-parameter families, by (iv).

\emph{Step 8 (clause (vi)).} Clause (iii) of Theorem \claimref{synthesis.one-loop-vacuum-polarisation-extends-without-a-regulator-and-ward-preserving-extensions-differ-only-by-the-renormalisation-constants}: the covariant, reflection-symmetric local differences of the pair that preserve the Ward-Takahashi identity form a space of dimension $2$ (the forced relation $c_{V} = -c_{2}$). That count is made in $\Cl(1,3)$, in the momenta, under the reflection $v \mapsto \gamma_{0}v\gamma_{0}$; it is carried here to position-space kernels in $\Cl(4,0)$ under rotations and reflections, where the covariant reflection-symmetric local terms of the allowed orders are $c_{m}\delta + c_{2}\nabla\delta$ and $c_{V}\,n\,\delta(u)\delta(v)$ (the ranks are not recomputed in $\Cl(4,0)$). The relation is rederived directly: $\operatorname{div}$ of $c_{V}\,n\,\delta\delta$ is $c_{V}(\nabla_{v} - \nabla_{u})\delta(u)\delta(v)$, and $\delta(v)(c_{m}\delta + c_{2}\nabla\delta)(u) - \delta(u)(c_{m}\delta + c_{2}\nabla\delta)(v) = c_{2}(\nabla_{u} - \nabla_{v})\delta(u)\delta(v)$, so the identity is preserved exactly when $c_{V} = -c_{2}$, with $c_{m}$ free. Together with one pair obeying the identity, clause (iii) here, this gives the stated family on the Euclidean section; nothing follows from it for the time-ordered products in signature $(1,3)$ without the Wick rotation of the vertex kernel, which is not stored. By Step 6 a change of $\ell$ at $m = 0$ moves along the second direction.

\emph{Status.} (i), (iv) and (v) are computed: algebra in $\Cl(4,0)$ and exact pairings, the passage from the test functions of the check to all test functions being the one-constant argument of Step 7. (ii) and (iii) are proved, given the imported extension theorem only for the statement about scaling degrees in Step 4 and given the standard properties of the massive Euclidean propagator (fundamental solution on both sides, exponential decay, local integrability, the stored expansion), which are not derived here. (vi) is the stored count, carried to the Euclidean section, plus (iii); it is a statement on the Euclidean section and does not discharge the imported existence for the time-ordered products. Standard physics throughout.
\end{proof}
